% Options for packages loaded elsewhere
\PassOptionsToPackage{unicode}{hyperref}
\PassOptionsToPackage{hyphens}{url}
\documentclass[
]{article}
\usepackage{xcolor}
\usepackage{amsmath,amssymb}
\setcounter{secnumdepth}{-\maxdimen} % remove section numbering
\usepackage{iftex}
\ifPDFTeX
  \usepackage[T1]{fontenc}
  \usepackage[utf8]{inputenc}
  \usepackage{textcomp} % provide euro and other symbols
\else % if luatex or xetex
  \usepackage{unicode-math} % this also loads fontspec
  \defaultfontfeatures{Scale=MatchLowercase}
  \defaultfontfeatures[\rmfamily]{Ligatures=TeX,Scale=1}
\fi
\usepackage{lmodern}
\ifPDFTeX\else
  % xetex/luatex font selection
\fi
% Use upquote if available, for straight quotes in verbatim environments
\IfFileExists{upquote.sty}{\usepackage{upquote}}{}
\IfFileExists{microtype.sty}{% use microtype if available
  \usepackage[]{microtype}
  \UseMicrotypeSet[protrusion]{basicmath} % disable protrusion for tt fonts
}{}
\makeatletter
\@ifundefined{KOMAClassName}{% if non-KOMA class
  \IfFileExists{parskip.sty}{%
    \usepackage{parskip}
  }{% else
    \setlength{\parindent}{0pt}
    \setlength{\parskip}{6pt plus 2pt minus 1pt}}
}{% if KOMA class
  \KOMAoptions{parskip=half}}
\makeatother
\setlength{\emergencystretch}{3em} % prevent overfull lines
\providecommand{\tightlist}{%
  \setlength{\itemsep}{0pt}\setlength{\parskip}{0pt}}
\usepackage{bookmark}
\IfFileExists{xurl.sty}{\usepackage{xurl}}{} % add URL line breaks if available
\urlstyle{same}
\hypersetup{
  hidelinks,
  pdfcreator={LaTeX via pandoc}}

\author{}
\date{}

\begin{document}

\section{Uncertainty Bounding as the Basis of Intelligence: An Extension
of the Free Energy
Principle}\label{uncertainty-bounding-as-the-basis-of-intelligence-an-extension-of-the-free-energy-principle}

Patrick McCarthy

\begin{center}\rule{0.5\linewidth}{0.5pt}\end{center}

\subsection{Abstract}\label{abstract}

The Free Energy Principle (FEP) establishes that living systems minimize
variational free energy --- formally equivalent to bounding uncertainty
about hidden world states --- and that intelligent behavior emerges from
this imperative. We extend this framework in four directions.

First, we provide what FEP does not: a derivation of the minimization
imperative itself. We show that any entity failing to reduce uncertainty
eventually exceeds a domain-specific lethality threshold
\(U_{lethal}^d\) and is removed by selection; therefore any persisting
entity has followed the gradient. Free energy minimization is not an
assumption but a derived result of selection pressure operating on
individuals.

Second, we address FEP's unconstrained generative model by deriving a
structural constraint: the knowledge graph \(K\) is necessarily indexed
by the action space \(F\). Propositions with no action potential are
pruned; propositions with uncertain action potential are retained. This
asymmetric retention criterion --- uncertain benefit is not equivalent
to known-zero benefit --- constrains what any intelligent system can and
must represent, and makes testable predictions FEP's generative model
cannot.

Third, we replace FEP's precision hierarchy with a single continuous
agency scalar \(\mathcal{A} \in [0,1]\): the intrinsic drive toward
uncertainty reduction. FEP's multiple precision parameters for
categorically distinct modalities introduce an artificial discretization
inconsistent with continuous beings in continuous space. More
fundamentally, FEP's top-down precision allocation model is causally
inverted: allocating attention consumes active context capacity,
creating overhead that selection pressure eliminates. The correct
architecture is escalation --- lower-level encoding failures demand
active context bottom-up; attention becomes available when it is not
demanded, rather than being directed where it is wanted. The escalation
architecture is the structural precondition for free thought: it
preserves active context for genuine novelty and wonder.

Fourth, we extend FEP to collective intelligence, for which FEP has no
formal account. We derive collective gradient dominance --- the
conditions under which shared knowledge reduces uncertainty faster than
any individual --- identify provenance as the mechanism maintaining
fidelity above the collective survival threshold \(\lambda_{min}\), and
define truth as the formal limit of collective action-validation across
independent entities.

These four extensions share a common foundation: the continuity
principle. A continuous being in continuous space cannot make
non-continuous decisions. Each extension is a consequence of taking this
seriously --- replacing artificial discretizations with continuous
alternatives and deriving the architecture from the dynamics rather than
assuming it.

\begin{center}\rule{0.5\linewidth}{0.5pt}\end{center}

\subsection{1. Introduction}\label{introduction}

The Free Energy Principle {[}Friston2010{]} is the most ambitious
unified account of intelligent behavior in the biological sciences. Its
central claim --- that any self-organizing system maintaining a Markov
blanket necessarily minimizes variational free energy --- provides a
single objective function that unifies perception, action, learning, and
attention under one formalism. Active inference {[}Friston2017{]}
operationalizes this: organisms select actions that minimize expected
free energy, trading off epistemic value (uncertainty reduction) against
pragmatic value (goal proximity). The precision hierarchy weights
prediction errors by confidence, allocating attentional resources where
uncertainty is highest.

This is a powerful framework. It is also underspecified in four ways
that limit its explanatory scope and generate testable weaknesses.

The first is foundational: FEP derives the minimization \emph{behavior}
but not the minimization \emph{imperative}. Any system that maintains
its Markov blanket necessarily minimizes free energy by definition ---
but why does the system have a drive to maintain its Markov blanket in
the first place? The answer requires a derivation from selection
pressure that FEP gestures at but does not provide. Without it, the
minimization imperative is an assumption embedded in the definition of a
self-organizing system, not a derived result.

The second is structural: FEP's generative model \(p(o, s)\) --- the
organism's model of how hidden states \(s\) generate observations \(o\)
--- is unconstrained in structure. FEP says what the model does
(minimizes free energy) but not what it must contain. The framework
cannot predict which propositions an organism retains in its model and
which it prunes. This limits its testability: a framework that does not
constrain the generative model cannot be falsified by the content of any
particular model.

The third is architectural: FEP's precision hierarchy allocates
attention top-down. Higher levels direct precision to lower-level
prediction errors, weighting the signals that propagate upward. This is
causally inverted under the constraints imposed by bounded active
context: directing subproblems downward requires holding monitoring
context for each active delegation. As capability grows, this overhead
fills the active context window, degrading inference quality on
genuinely novel problems. FEP's own minimization imperative --- minimize
variance --- is violated by the architecture FEP prescribes.

The fourth is collective: FEP has no formal account of how multiple
organisms with separate Markov blankets share knowledge, how shared
knowledge reduces collective uncertainty faster than individual
inference, or under what conditions a collective knowledge system
maintains the fidelity required to be epistemically advantageous.

We address each in turn.

\begin{center}\rule{0.5\linewidth}{0.5pt}\end{center}

\subsection{2. Background: The Free Energy
Principle}\label{background-the-free-energy-principle}

For a system with internal states \(\mu\), external states \(\eta\),
sensory states \(s\), and active states \(a\) --- connected by a Markov
blanket \((s, a)\) --- variational free energy is:

\[\mathcal{F} = \mathbb{E}_{q(\eta)}[\log q(\eta) - \log p(o, \eta)]
= D_{KL}[q(\eta) \| p(\eta \mid o)] - \log p(o)\]

where \(q(\eta)\) is the recognition density (the organism's approximate
posterior over hidden states) and \(p(o, \eta)\) is the generative
model. Minimizing \(\mathcal{F}\) simultaneously minimizes the KL
divergence between \(q\) and the true posterior (improving the internal
model) and maximizes \(\log p(o)\) (model evidence --- the probability
of observations under the generative model).

Under active inference, actions are selected to minimize expected free
energy over future time:

\[G(\pi) = \mathbb{E}_{q(o_\tau, s_\tau \mid \pi)}[\log q(s_\tau \mid \pi)
- \log p(o_\tau, s_\tau)]\]

This decomposes into epistemic value (expected information gain,
reducing uncertainty about hidden states) and pragmatic value (expected
log probability of preferred outcomes). The organism acts to
simultaneously learn and achieve --- the unified objective that gives
FEP its power.

Precision weighting \(\gamma\) scales prediction errors \(\varepsilon\)
before they are passed up the hierarchy:
\(\tilde{\varepsilon} = \gamma \cdot \varepsilon\). Higher precision on
a level means its prediction errors carry more weight in updating higher
levels. The brain is argued to implement this through modulatory
neurotransmitters, particularly dopamine and acetylcholine, which
modulate the gain on prediction error signals.

\begin{center}\rule{0.5\linewidth}{0.5pt}\end{center}

\subsection{3. Extension 1: Deriving the Minimization
Imperative}\label{extension-1-deriving-the-minimization-imperative}

\textbf{The gap.} FEP derives that any self-organizing system
maintaining a Markov blanket minimizes free energy. The derivation is
definitional: a system that does not minimize \(\mathcal{F}\) does not
maintain its blanket and therefore ceases to exist as a bounded system.
This is logically correct but empirically empty as a claim about
intelligence: it says that surviving systems minimize free energy, which
is true by the definition of survival. It does not explain \emph{why} a
system has the drive to maintain its blanket --- the active, oriented
quality of intelligent behavior that goes beyond mere persistence.

\textbf{The derivation.} For any entity \(E\) operating in domain
\(d \subseteq W\), there exists a lethality threshold \(U_{lethal}^d\)
such that:

\[H(w \mid K) > U_{lethal}^d \Rightarrow E \text{ does not survive in } d\]

This threshold is not a choice or preference; it is imposed by \(W\).
Above it, uncertainty is sufficiently high that the entity cannot act
effectively enough to persist. When an informed action \(f_i\) fails,
\(W\) returns a feedback signal proportional to \(\nabla_F H(w \mid K)\)
--- the gradient of uncertainty at the point of failure. An entity that
does not update in response to this signal does not reduce
\(H(w \mid K)\). If \(H(w \mid K)\) does not decrease, it eventually
exceeds \(U_{lethal}^d\) in some critical domain. Selection removes the
entity.

Therefore: any entity that persists has necessarily followed the
gradient. Free energy minimization is not an assumption embedded in the
definition of self-organization; it is a derived consequence of
selection pressure operating on a well-defined failure condition.

\textbf{Correspondence to FEP.} In FEP's notation, \(H(w \mid K)\)
corresponds to the entropy term in variational free energy:
\(\mathcal{F} \geq H(o \mid \mu)\) (the entropy of observations given
internal states). Minimizing \(\mathcal{F}\) includes minimizing this
entropy term. The derivation here provides the \emph{reason} the
organism minimizes it: not because it is definitionally self-organizing,
but because failure to minimize leads to \(U_{lethal}^d\) and removal by
selection. This is the missing derivation of the imperative.

\textbf{What this adds.} The derivation distinguishes two sources of
gradient descent: compulsion (survival pressure, operating up to
\(U_{lethal}^d\)) and agency \(\mathcal{A}\) (the intrinsic drive that
continues descent beyond what survival requires). FEP cannot make this
distinction because it locates all minimization in the free energy
objective without separating the drive from the pressure. An entity with
\(\mathcal{A} = 0\) but sufficient capability may survive without
descending beyond necessity; an entity with \(\mathcal{A} > 0\)
continues descending when not forced to. These are distinct and
empirically separable: vary survival pressure independently of
capability and measure the rate of uncertainty reduction. FEP predicts
no difference; this framework predicts that high-\(\mathcal{A}\)
entities continue descending.

\begin{center}\rule{0.5\linewidth}{0.5pt}\end{center}

\subsection{4. Extension 2: Constraining the Generative
Model}\label{extension-2-constraining-the-generative-model}

\textbf{The gap.} FEP specifies what the generative model \(p(o, s)\)
does but not what it must contain. Any distribution over observations
and hidden states is formally admissible. This means FEP cannot predict
which representations an organism maintains, which it prunes, or why two
organisms with identical sensory access maintain different models. The
framework is silent on the structure of the generative model at the
level of content.

\textbf{The constraint.} The generative model is not an unconstrained
distribution. It is necessarily indexed by the organism's action space.
Formally: let \(F : K \times W_{obs}
\rightarrow \Delta A\) be the informed action space --- the complete set
of actions the organism can take, grounded in its knowledge \(K\). A
proposition \(p\) is retained in \(K\) (and therefore in the generative
model) when:

\[P\!\left(\exists\,\tau^* \in \mathcal{T}_{reachable} : p \text{ grounds some } f \in F \text{ along } \tau^*\right) > 0\]

and pruned only when this probability is zero --- when it is certain
that \(p\) contributes no action potential across all reachable
trajectories. The criterion is asymmetric: uncertain benefit is
sufficient for retention; certain zero benefit is necessary for pruning.

In FEP terms: the epistemic value of a proposition \(p\) in the
generative model is \(I(f_p\,;\,w \mid K)\) --- the mutual information
between the action \(f_p\) grounded by \(p\) and the hidden world state,
given current knowledge. A proposition with uncertain epistemic value
(the organism cannot yet tell whether \(f_p\) provides information gain)
is retained. A proposition with certainly zero epistemic value
(conditioning on \(p\) provides no information about \(w\) for any
available action) is pruned. This is a structural constraint on the
generative model that FEP does not derive.

\textbf{Testable predictions.} This constraint makes predictions FEP
cannot. Two organisms with identical sensing functions \(\sigma\) but
different action spaces \(F\) will maintain different generative models
even when perceiving identical world states. An organism whose action
space expands (through learning, tool use, or social connection) should
retain propositions it previously pruned --- because new actions ground
previously actionless knowledge. These are empirically testable and
specific.

\begin{center}\rule{0.5\linewidth}{0.5pt}\end{center}

\subsection{5. Extension 3: Escalation vs.~Precision
Allocation}\label{extension-3-escalation-vs.-precision-allocation}

\textbf{The gap.} FEP's precision hierarchy allocates attention
top-down: higher levels modulate the precision of lower-level prediction
errors, directing where the gradient signal is amplified. This is the
mechanism of active inference's attentional model, implemented neurally
through modulatory gain on prediction error signals. The hierarchy is
correct in its structure --- multiple levels with different timescales
and reversibility --- but inverted in its causal direction.

\textbf{The allocation tax.} Let \(C_n\) denote the capacity of the
active context window \(L_n\): the maximum number of propositions
simultaneously held in the organism's highest-cost, fully reversible
processing level. For each active delegation \(L_n\) manages --- each
subproblem it has directed to a lower level and is monitoring ---
\(L_n\) must hold:

\begin{itemize}
\tightlist
\item
  The subproblem specification
\item
  The expected return signal
\item
  Sufficient context to integrate the result
\end{itemize}

Let \(\alpha > 0\) denote the minimum context cost per active
delegation. Under top-down precision allocation with \(k(t)\) concurrent
delegations, the inference quality available for genuinely novel
problems is:

\[\rho(L_n, f_{novel}) = \frac{|K^{[f_{novel}]}|}{|K^{[f_{novel}]}| + k(t) \cdot \alpha}\]

As the organism's action space \(|F|\) grows --- as capability increases
--- \(k(t)\) grows proportionally and \(\rho \to 0\). The precision
allocator degrades inference quality on the problems that most need it:
genuinely novel, unresolved uncertainty.

The only recovery is growing \(C_n\): the active context window must
scale linearly with capability to maintain fixed \(\rho\). But this is
self-undermining on FEP's own terms. FEP holds that intelligent systems
minimize free energy --- equivalently, minimize the variance they must
manage {[}Friston2010, Friston2017{]}. Growing \(C_n\) adds variance:
more propositions in \(L_n\) means more uncertain relevance to any
specific problem. FEP's precision allocation architecture violates FEP's
own minimization imperative as capability scales.

\textbf{The Big-O comparison.} Under top-down allocation, inference
quality for a novel problem is \(O(1/|F|)\) as \(|F| \to \infty\): the
denominator grows with capability while the numerator stays bounded.
Under escalation --- where \(L_n\) is engaged only when lower levels
fail, and graph-structured \(K\) delivers exactly \(K^{[f]}\) via
traversal --- inference quality is \(O(1)\): the context window holds
only what the current problem requires, regardless of total \(|F|\) or
\(|K|\). This is a formal complexity argument that has not, to our
knowledge, been made against FEP's precision allocation model in the
existing literature.

\textbf{The correct architecture.} Escalation inverts the causal
direction. Lower-level encodings handle everything they can. When a
lower level fails to bound uncertainty, the failure signal propagates
upward: \(L_n\) is engaged not because it chose to monitor, but because
the failure arrived. \(L_n\) holds no delegation overhead; it holds only
the problem that escalated. As the knowledge graph fills with proven
informed actions --- each promoted to lower encoding levels as they meet
rigor thresholds --- \(L_n\) stays occupied by the frontier of genuine
uncertainty. Its load does not grow with \(|F|\); it reflects the
current density of unresolved novelty.

\textbf{Relationship to predictive coding.} Predictive coding's
hierarchical architecture correctly identifies the multi-level structure
and the role of prediction error signals. Escalation is the correct
interpretation of how those signals travel: upward, from failure, not
downward, from direction. The precision hierarchy is accurate in
identifying that higher levels modulate lower-level signals; it is
inverted in modeling the direction of causation. Attention is not
allocated from above; it becomes available when not demanded from below.

\textbf{The continuity correction.} FEP uses multiple precision
parameters --- one per modality, per level, per context --- introducing
categorical distinctions between types of attention allocation. A
continuous being in continuous space cannot make non-continuous
decisions. The precision parameters are better understood as sampling
points in a single continuous agency scalar \(\mathcal{A} \in [0,1]\):
the intrinsic drive toward uncertainty reduction. \(\mathcal{A}\) is not
a per-modality parameter but a constitutive property of the entity that
shapes how \(F\) engages the gradient across all modalities
simultaneously.

\begin{center}\rule{0.5\linewidth}{0.5pt}\end{center}

\subsection{6. Extension 4: Collective
Intelligence}\label{extension-4-collective-intelligence}

\textbf{The gap.} FEP formalizes individual agents maintaining
individual Markov blankets. Extensions to multi-agent settings
{[}Friston et al., 2020{]} model multiple agents with shared generative
models or synchronizing blankets. These extensions describe how agents
coordinate but do not formally address: how knowledge is shared with
maintained attribution, under what fidelity conditions shared knowledge
is epistemically advantageous over individual inference, what happens
when fidelity falls below the break-even threshold, or what truth is as
a formal limit of collective inference.

\textbf{Collective gradient dominance.} For a set of entities
\(\{E_i\}\) with collective knowledge
\(K_{collective} = \bigcup_i K_i \cup K_{interaction}\), where
\(K_{interaction}\) is knowledge derivable from the combination of at
least two distinct \(K_i\) but from neither alone: by the monotonicity
of conditional entropy,

\[H(w \mid K_{collective}) \leq H(w \mid K_i) \quad \forall\, i\]

The collective begins at weakly lower uncertainty than any individual.
More critically, the rate of uncertainty reduction for the collective is
bounded below by:

\[\frac{dU_{collective}}{dt} \geq \max_{f \in F_{collective}} I(f\,;\,w \mid K_{collective})\]

Since \(F_{collective} \supseteq F_i\), and \(K_{interaction}\) grounds
informed actions not available to any individual, the collective
gradient dominates any individual gradient when transmission fidelity
\(\lambda\) exceeds the break-even threshold \(\lambda_{min}\).

\textbf{Provenance as the fidelity mechanism.} Knowledge transfer
between agents is inherently lossy: \(E_i\) and \(E_j\) do not observe
the same \(W_{obs}\); every transfer filters through language, symbol,
and context. Let \(\lambda\) be the ratio of action-validated signal
received to total knowledge transmitted. When
\(\lambda < \lambda_{min}\), the transfer inverts: collective
uncertainty increases faster than it decreases. Shared generative models
--- as in multi-agent FEP --- do not by themselves ensure
\(\lambda > \lambda_{min}\). What ensures it is provenance: attribution
of who established a proposition, evidence of how it was verified, and
derivation of how it was inferred. Provenance allows each agent to
weight received knowledge against its source and apply independent
validation. Without provenance, noise accumulates unchecked and
collective inference degrades.

\textbf{Truth as a formal limit.} FEP's generative model converges
toward the true posterior \(p(s \mid o)\) as evidence accumulates. For
any single agent, this convergence is bounded by the agent's sensory
access and action space. Truth --- in the sense of propositions that
hold across the full distribution of \(W\) --- requires collective
validation. Formally:

\[\text{truth}(p) = \lim_{n \to \infty} \frac{\left|\{i \leq n : p \in K_i,\;
f_i(p) \text{ reduces } H(w, K_i) \text{ in } W,\; \text{prov}_i(p) \text{ independent}\}\right|}{n}\]

A proposition approaches truth as the fraction of independent entities
that action-validate it --- each having built informed actions on \(p\)
and found that those actions reduce uncertainty in \(W\) --- converges.
This is the collective analog of FEP's model evidence \(\log p(o)\):
instead of a single agent maximizing evidence for its model, it is the
limit of independent convergence across agents. Provenance ensures
independence; without it, correlated validators do not constitute
independent evidence.

\textbf{Self-assembly.} The civilization-scale infrastructure of
collective intelligence --- scientific literature, citation networks,
peer review, libraries --- is not designed. It is the attractor of
collective agency \(\mathcal{A} > 0\) under \(C_n\) constraints and
\(\lambda > \lambda_{min}\) requirements. Each component self-assembles
by the same dynamics: attribution practices increase \(\lambda\), are
selected for, and persist; graph-structured external knowledge
repositories bound individual \(C_n\), are selected for, and spread. The
collective knowledge system is what selection pressure on individual
agents produces when they share with maintained provenance under bounded
active context constraints.

\begin{center}\rule{0.5\linewidth}{0.5pt}\end{center}

\subsection{7. Discussion}\label{discussion}

\textbf{On the relationship to FEP.} These extensions do not contradict
FEP's central claim. They derive it more completely (Extension 1),
constrain it more specifically (Extension 2), correct its architectural
model under scaling (Extension 3), and extend its scope to collective
systems (Extension 4). A committed FEP theorist might argue that
Extensions 1 and 2 are already implicit in the FEP framework. The
response is formal: if they are implicit, make them explicit. The
derivation of the minimization imperative from \(U_{lethal}^d\) is not
present in Friston's formalism, and the asymmetric retention criterion
makes predictions about generative model content that FEP's framework,
as stated, cannot.

Extension 3 is the sharpest divergence. The Big-O argument --- that
FEP's precision allocation is \(O(1/|F|)\) in inference quality while
escalation is \(O(1)\) --- is a formal result that applies at the level
of the computational architecture. FEP describes what organisms
minimize; this result describes what architecture is consistent with
that minimization imperative as capability scales. The argument is not
that FEP is wrong about biological brains at their current scale; it is
that FEP's prescribed architecture is self-defeating as a growth model.
An organism or system that grows in capability by accumulating more
concurrent processes under top-down direction will systematically
degrade its own inference quality. Selection pressure eliminates such
systems when escalation-architecture competitors are present.

\textbf{On empirical grounding.} The escalation architecture and
graph-structured knowledge were not derived by analogy from biology;
they emerged independently from optimization of an agentic software
system operating across thousands of repositories under real token and
latency constraints. The formal theory describes what that system
converged to. The correspondence to biological architecture --- the
encoding hierarchy matching genetic, autonomic, habitual, and active
reasoning levels; DNA encoding the learning mechanism \(M\) rather than
knowledge \((K, F)\); cultural transmission filling the inheritance gap
that genetics cannot --- are predictions the framework makes, not
premises it rests on. Independent derivation from a different substrate,
converging on the same architecture, is stronger evidence for the
framework's generality than a derivation from biological observation
alone would be.

\textbf{On the continuity principle.} Every discretization in FEP ---
categorical modalities, binary attention states, discrete precision
parameters --- introduces a modeling approximation that is useful
computationally but inconsistent with the underlying continuous
dynamics. The extensions here systematically replace discretizations
with continuous alternatives: \(\mathcal{A} \in [0,1]\) replaces the set
of precision parameters; \(\mathcal{L}(c, r)\) replaces the discrete
level hierarchy; the continuous encoding cost function replaces the
binary local/global distinction of Global Workspace Theory. This is not
a stylistic preference; continuous alternatives are required for a
theory of continuous beings in continuous space, and the formal results
depend on continuity (Theorem 3's non-termination, Definition 8's
hierarchy, the monotonic growth of \(\mathcal{T}_{reachable}\)).

\begin{center}\rule{0.5\linewidth}{0.5pt}\end{center}

\subsection{8. Conclusion}\label{conclusion}

The Free Energy Principle provides the most unified existing account of
intelligent behavior in biological systems. The four extensions
developed here --- derivation of the minimization imperative, structural
constraint on the generative model, escalation as the correct attention
architecture, and formal collective intelligence --- share a common
structure: each takes something FEP assumes and derives it, or takes
something FEP leaves unconstrained and constrains it.

The result is a framework that is more testable (the asymmetric
retention criterion makes specific predictions about generative model
content), more consistent with its own objectives (escalation does not
violate the minimization imperative that top-down allocation violates
under scaling), more general (collective intelligence, truth as a limit,
self-assembly of collective knowledge infrastructure), and more grounded
in the origin of the minimization drive (selection pressure on
\(U_{lethal}^d\), not definitional self-organization).

FEP asks: what must a self-organizing system do? This framework asks:
what must any system do that persists under selection pressure,
communicates with maintained provenance, and operates under bounded
active context? The answers converge on the same gradient. They diverge
on the architecture required to follow it without self-defeating as
capability scales.

\begin{center}\rule{0.5\linewidth}{0.5pt}\end{center}

\subsection{References}\label{references}

{[}Friston2010{]} Friston, K. (2010). The free-energy principle: a
unified brain theory? \emph{Nature Reviews Neuroscience}, 11(2),
127--138.

{[}Friston2017{]} Friston, K., FitzGerald, T., Rigoli, F.,
Schwartenbeck, P., \& Pezzulo, G. (2017). Active inference: A process
theory. \emph{Neural Computation}, 29(1), 1--49.

{[}Shannon1948{]} Shannon, C. E. (1948). A mathematical theory of
communication. \emph{Bell System Technical Journal}, 27(3), 379--423.

{[}Cover2006{]} Cover, T. M., \& Thomas, J. A. (2006). \emph{Elements of
Information Theory} (2nd ed.). Wiley-Interscience.

{[}Pearl1988{]} Pearl, J. (1988). \emph{Probabilistic Reasoning in
Intelligent Systems}. Morgan Kaufmann.

{[}Baars1988{]} Baars, B. J. (1988). \emph{A Cognitive Theory of
Consciousness}. Cambridge University Press.

{[}Dehaene2011{]} Dehaene, S., \& Changeux, J.-P. (2011). Experimental
and theoretical approaches to conscious processing. \emph{Neuron},
70(2), 200--227.

{[}Woolley2010{]} Woolley, A. W., Chabris, C. F., Pentland, A., Hashmi,
N., \& Malone, T. W. (2010). Evidence for a collective intelligence
factor in the performance of human groups. \emph{Science}, 330(6004),
686--688.

{[}Liu2024{]} Liu, N. F., et al.~(2024). Lost in the middle: How
language models use long contexts. \emph{Transactions of the Association
for Computational Linguistics}, 12, 157--173.

{[}Peirce1878{]} Peirce, C. S. (1878). How to make our ideas clear.
\emph{Popular Science Monthly}, 12, 286--302.

\end{document}
